% TEMPLATE-MILC-v3.tex - plantilla maestra UNICA de las guias MILC v3.
%
% La usan las cuatro familias de guias, cada una con su driver:
%   · Grados 6-11          -> scripts/build-guias-g11.py
%   · Bebras               -> scripts/build-guias-bebras.py
%   · Semillero            -> scripts/build-guias-semillero.py
%   · Territorio interior  -> scripts/build-guias-territorio-interior.py
%
% Antes habia tres copias casi identicas (~400 lineas iguales, 6 distintas):
% cada mejora habia que aplicarla tres veces, y en la practica no se hacia
% (el motor de recursos vivia solo en dos de ellas). Lo que varia entre
% familias esta parametrizado en 6 placeholders que cada driver llena
% (se nombran aqui SIN su sintaxis real: el builder reemplaza por texto plano,
% tambien dentro de los comentarios, y un valor multilinea rompe el preambulo):
%
%   HEADER_IZQ        encabezado izquierdo de pagina
%   HEADER_DER        encabezado derecho de pagina
%   PORTADA_ETIQUETA  rotulo sobre el numero grande (GRADO/MOMENTO/MODULO)
%   PORTADA_SERIE     linea de serie de la portada
%   PORTADA_ID        identificador de la guia en la portada
%   PORTADA_CIRCULO   el circulito de grado (°), o vacio si no aplica
%   PDF_TITULO        titulo en texto plano (sin \textbf ni comillas TeX) para
%                     los metadatos del PDF (/Title); lo produce lib_xelatex.texto_pdf
%
% Composicion (auditoria 2026-09-03): las bandas de estacion piden espacio
% (\Needspace*) y prohiben el salto justo despues (\nopagebreak), asi no quedan
% huerfanas al pie; las cajas largas (softbox, drawbox, infoband, puentebox) son
% `breakable`, asi una caja de media pagina ya no arrastra todo a la siguiente
% dejando la anterior al 40 %. Todo texto sobre mostaza o verde va en milcNegro
% (blanco sobre #E5B400 da 1,93:1 y sobre #58A923 2,95:1; ninguno pasa WCAG AA).
% El resto de claves las documenta content/guias/_SCHEMA.md. El script valida
% que no queden placeholders sin reemplazar antes de escribir el archivo final.

% Etiquetado PDF (latex-lab): /Lang, estructura y texto alternativo de las
% figuras, para lectores de pantalla. Debe ir ANTES de \documentclass.
\DocumentMetadata{lang=es-CO,tagging=on}
\documentclass[11pt,letterpaper]{article}
% headheight=24pt: la cabecera derecha lleva dos lineas (identidad de la guia
% y estacion actual). headsep se acorta en la misma medida para que el bloque
% de texto quede exactamente donde estaba con headheight=12pt.
\usepackage[margin=2.0cm,headheight=24pt,headsep=13pt]{geometry}
\usepackage[table]{xcolor}
\usepackage{fontspec}
\usepackage[spanish]{babel}
\usepackage{tikz}
% Librerías TikZ para los diagramas de las guías (bloque `recursos:`):
%  · babel  → babel en español vuelve activos caracteres como «>», y TikZ los usa
%             en sus opciones (>=stealth, ->). Sin esta librería, cualquier
%             diagrama con flechas revienta con «\language@active@arg> has an
%             extra }». Debe ir ANTES de que se dibuje nada.
%  · positioning        → sintaxis «below left=2pt and 3pt».
%  · decorations.pathreplacing → llaves ({brace}) para señalar rangos.
\usetikzlibrary{babel,positioning,decorations.pathreplacing}
\usepackage{graphicx}
% Circuitos: el trabajo de electrónica se hace hoy con micro:bit y sus sensores
% integrados (MakeCode, sin cableado), así que no se necesitan esquemáticos.
% Si algún día entran montajes con protoboard, basta descomentar esta línea:
% los diagramas .tex/.tikz declarados en `recursos:` ya se incrustan solos.
% \usepackage{circuitikz}
\usepackage[most]{tcolorbox}
\usepackage{tabularx}
\usepackage{array}
\usepackage{fancyhdr}
\usepackage{titlesec}
\usepackage[document]{ragged2e}  % bandera a la derecha en todo el cuerpo
\usepackage{enumitem}
\usepackage{microtype}
\usepackage{etoolbox}
\usepackage{needspace}
% hyperref al final del preambulo: metadatos del PDF (titulo, autor, idioma,
% familia), marcadores por estacion (\pdfbookmark) y enlaces sin recuadro.
\usepackage[hidelinks,
  pdftitle={Cosecha P1 — manifiesto del técnico crítico},
  pdfauthor={PhD. Álvaro Cárdenas Orozco},
  pdfsubject={Tecnología e Informática · Grado 9 · Guía 10 · MILC v3},
  pdflang={es-CO},
  bookmarksopen=true,bookmarksnumbered=false]{hyperref}

% Helvetica Neue no trae la flecha U+2192, asi que cada «→» escrito literalmente
% en el YAML se compilaba a NADA, en silencio (462 flechas en 61 guias: rutas del
% tipo «paleta → tipografia → lockup» salian rotas). Mapeamos el caracter a su
% version matematica, que si tiene glifo. Asi el autor puede seguir escribiendo
% «→» comodamente en el YAML.
\usepackage{newunicodechar}
\newunicodechar{→}{\ensuremath{\rightarrow}}
% Mismo problema con los simbolos de marca y vinetas: el «✓» de «una palomita ✓»
% salia como cuadrito vacio en el PDF de todas las guias que lo usaban.
\usepackage{pifont}
\newunicodechar{✓}{\ding{51}}
\newunicodechar{✗}{\ding{55}}
\newunicodechar{✔}{\ding{52}}
\newunicodechar{•}{\textbullet}
\newunicodechar{≥}{\ensuremath{\geq}}
\newunicodechar{≤}{\ensuremath{\leq}}

\setmainfont{Helvetica Neue}
\setsansfont{Helvetica Neue}
\setlength{\parindent}{0pt}
\setlength{\parskip}{6pt}
% Sin particion de palabras: en 6.º-7.º una palabra cortada por guion al final
% de linea («Comuni-cacion») frena la lectura mucho mas que una linea corta.
\hyphenpenalty=10000
\exhyphenpenalty=10000
\pretolerance=2000
\tolerance=3000
\setlength{\emergencystretch}{3em}
\linespread{1.10}
\renewcommand{\arraystretch}{1.25}
\setlist{leftmargin=5mm,itemsep=1mm,topsep=1mm}
\newcolumntype{Y}{>{\RaggedRight\arraybackslash}X}

\definecolor{milcMagenta}{HTML}{E6007E}
\definecolor{milcVino}{HTML}{5A0038}
\definecolor{milcTurquesa}{HTML}{0093A5}
\definecolor{milcVerde}{HTML}{58A923}
\definecolor{milcMostaza}{HTML}{E5B400}
\definecolor{milcOcre}{HTML}{B86600}
\definecolor{milcNegro}{HTML}{241816}
\definecolor{milcGris}{HTML}{F7F7F4}
\definecolor{milcLinea}{HTML}{E8E2DD}
\definecolor{milcGradePrimary}{HTML}{7C3AED}
\definecolor{milcGradeDark}{HTML}{4C1D95}
\definecolor{milcGradeSoft}{HTML}{EDE9FE}
\definecolor{milcGradeAccent}{HTML}{CFE7FF}
\definecolor{milcGradeText}{HTML}{FFFFFF}
\color{milcNegro}

\pagestyle{fancy}
\fancyhf{}
% Cabecera de dos lineas: identidad de la guia arriba y, a la derecha y en
% gris, la estacion en curso (\markright desde \milcestacion). Antes de la
% primera estacion la segunda linea va vacia; el \strut mantiene la altura
% para que la linea de arriba no baile entre paginas.
\lhead{\small Tecnología e Informática · Grado 9\\[-1pt]{\footnotesize\strut}}
\rhead{\small Guía 10 · MILC v3\\[-1pt]{\footnotesize\strut\color{milcNegro!65}\rightmark}}
\cfoot{\small\thepage}
\renewcommand{\headrulewidth}{0pt}

% Sobre mostaza (#E5B400) y verde (#58A923) el texto va en milcNegro; sobre el
% resto de la paleta, en blanco. Se usa dentro de fonttitle/fontupper, donde un
% \color posterior manda sobre coltitle.
\newcommand{\milctintasobre}[1]{%
  \ifboolexpr{test{\ifstrequal{#1}{milcMostaza}} or test{\ifstrequal{#1}{milcVerde}}}%
    {\color{milcNegro}}{\color{white}}}

\newtcolorbox{titlebox}[1]{
  enhanced,colback=#1,colframe=#1,arc=7mm,boxrule=0pt,
  left=7mm,right=7mm,top=3mm,bottom=3mm,
  before skip=8pt,after skip=8pt,
  fontupper=\sffamily\bfseries\Large\color{white}
}

\newtcolorbox{softbox}[3]{
  enhanced,breakable,pad at break=3mm,
  colback=#2,colframe=#1,borderline west={4pt}{0pt}{#1},
  arc=2mm,boxrule=.4pt,left=4mm,right=4mm,top=3mm,bottom=3mm,
  before skip=6pt,after skip=8pt,title={#3},coltitle=white,
  colbacktitle=#1,fonttitle=\sffamily\bfseries\milctintasobre{#1},fontupper=\RaggedRight,
  attach boxed title to top left={xshift=3mm,yshift=-2mm},
  boxed title style={arc=2mm, boxrule=0pt}
}

\newtcolorbox{drawbox}[2]{
  enhanced,breakable,pad at break=4mm,
  colback=white,colframe=#1,arc=4mm,boxrule=1pt,
  left=5mm,right=5mm,top=4mm,bottom=4mm,before skip=8pt,after skip=8pt,
  title={#2},coltitle=white,colbacktitle=#1,fonttitle=\sffamily\bfseries\milctintasobre{#1},
  fontupper=\RaggedRight,
  attach boxed title to top left={xshift=4mm,yshift=-2mm},
  boxed title style={arc=3mm, boxrule=0pt}
}

\newtcolorbox{infoband}[1]{
  enhanced,breakable,pad at break=2mm,colback=#1!8,colframe=#1!50,arc=2mm,boxrule=.5pt,
  left=4mm,right=4mm,top=2mm,bottom=2mm,before skip=4pt,after skip=4pt,
  fontupper=\small\RaggedRight
}

% Puente narrativo entre secciones (contrato editorial Conéctate).
% Da continuidad: "Acabas de X. Ahora vas a Y porque Z."
% Visualmente: banda izquierda en color del grado, texto en itálica pequeña.
\newtcolorbox{puentebox}{
  enhanced,breakable,pad at break=2mm,colback=milcGradeSoft,colframe=milcGradeDark,
  borderline west={3pt}{0pt}{milcGradeDark},
  arc=0mm,boxrule=0pt,left=5mm,right=4mm,top=2.5mm,bottom=2.5mm,
  before skip=4pt,after skip=8pt,
  fontupper=\itshape\small\color{milcNegro}\RaggedRight
}
% La flecha va en modo matematico a proposito: Helvetica Neue NO tiene el glifo
% U+2192, asi que \textrightarrow se compilaba a NADA («Missing character: There
% is no → in font Helvetica Neue») y el puente salia sin flecha. En modo
% matematico la flecha viene de la fuente matematica y si se dibuja.
\newcommand{\puente}[1]{%
  \begin{puentebox}%
    \textcolor{milcGradeDark}{$\rightarrow$}\hspace{2mm}#1%
  \end{puentebox}%
}

\newcommand{\linead}[1]{\noindent\color{milcLinea}\rule{#1}{0.5pt}\color{milcNegro}}
\newcommand{\respuesta}[1]{\par\vspace{2mm}\foreach \n in {1,...,#1}{\noindent\color{milcLinea}\rule{\linewidth}{0.5pt}\par\vspace{5mm}}\color{milcNegro}}
\newcommand{\checkbox}{\textcolor{milcGradeDark}{\fbox{\phantom{x}}}\hspace{5pt}}

% ─── Listas aireadas para las guías socioemocionales ────────────────────
% Componer las verificaciones como «\checkbox texto\\» deja el texto que salta
% de línea debajo del cuadrito y apelmaza el bloque. Estos entornos alinean la
% continuación y airean los ítems. Los usa Territorio Interior; cualquier
% familia puede hacerlo.
\newlist{checklista}{itemize}{1}
\setlist[checklista]{
  label=\textcolor{milcGradeDark}{\fbox{\phantom{x}}},
  leftmargin=9mm, labelsep=3.2mm, itemsep=2.4mm,
  topsep=2.5mm, parsep=0pt, partopsep=0pt, before=\RaggedRight
}
\newlist{pasoslista}{enumerate}{1}
\setlist[pasoslista]{
  label=\textcolor{milcGradeDark}{\sffamily\bfseries\arabic*.},
  leftmargin=8mm, labelsep=3mm, itemsep=2.8mm,
  topsep=1mm, parsep=0pt, partopsep=0pt, before=\RaggedRight
}

\newcommand{\phasecard}[4]{
\begin{tcolorbox}[enhanced,colback=#3,colframe=#2,arc=3mm,boxrule=.5pt,height=3.05cm,valign=center,halign=center,left=1.5mm,right=1.5mm,top=2mm,bottom=2mm]
{\sffamily\bfseries\fontsize{22}{24}\selectfont\textcolor{#2}{#1}}\par
{\sffamily\bfseries\small #4}
\end{tcolorbox}
}

% ── Piezas del rediseno 2026 (legibilidad 6.º-11.º) ───────────────────
% Banda de estacion: reemplaza el titlebox macizo. Titulo a la izquierda,
% dato practico (tiempo/modalidad) a la derecha, en una sola linea.
% Nunca huerfana: pide 9 lineas libres (o salta de pagina rellenando con
% \vfill) y prohibe el salto justo despues de la banda. Nueve y no seis:
% la caja partible que sigue necesita ~80pt (titulo adosado, relleno y dos
% lineas) para arrancar en la misma pagina; con menos, tcolorbox la manda
% entera a la siguiente y la banda se queda sola. El penalty 10000 va
% en `after`, ANTES del espacio posterior: un `after skip` (aunque sea 0pt)
% es un \vskip tras la caja, y ese glue es un punto de corte legal que un
% \nopagebreak escrito despues de \end{tcolorbox} ya no cierra. Cada banda
% deja marcador PDF y actualiza la cabecera derecha.
\newcounter{milcestacion}
\newcommand{\milcestacion}[2]{%
\Needspace*{9\baselineskip}%
\stepcounter{milcestacion}%
\pdfbookmark[1]{#2}{milcestacion\themilcestacion}%
\markright{#2}%
\begin{tcolorbox}[enhanced,colback=#1,colframe=#1,arc=2mm,boxrule=0pt,
  left=5mm,right=5mm,top=2.6mm,bottom=2.6mm,before skip=11pt,
  after={\par\nopagebreak\vskip7pt}]
{\sffamily\bfseries\fontsize{15}{18}\selectfont\milctintasobre{#1}#2}
\end{tcolorbox}}

% Tarjeta de paso numerada: sustituye las tablas «Pilar / Aplicacion», que
% eran formato de consulta docente, no de estudiante. El numero va en un
% circulo sobre el borde para que la secuencia se lea de un vistazo.
\newcommand{\milcpaso}[3]{%
\Needspace{4\baselineskip}%
\begin{tcolorbox}[enhanced,colback=white,colframe=milcLinea,arc=1.5mm,boxrule=.9pt,
  left=14mm,right=4mm,top=3mm,bottom=3mm,before skip=4pt,after skip=4pt,
  fontupper=\RaggedRight,
  overlay={\node[anchor=north west,circle,fill=#2,text=white,inner sep=0pt,
    minimum size=7.4mm,font=\sffamily\bfseries\small]
    at ([xshift=3.6mm,yshift=-3.2mm]frame.north west) {#1};}]
#3
\end{tcolorbox}}

% Casilla marcable de tamano util con lapiz (la anterior era \fbox{\phantom{x}},
% de alto variable segun el texto que la seguia).
\newcommand{\milcmarca}{\framebox[3.8mm]{\rule{0pt}{3.4mm}}}

% Fila marcable de lista suelta (checks de la estacion 1).
\newcommand{\milccheckrow}[1]{%
\par\vspace{1.8mm}\noindent
\makebox[6.4mm][l]{\raisebox{-.55mm}{\textcolor{milcNegro}{\milcmarca}}}%
\parbox[t]{\dimexpr\linewidth-6.4mm\relax}{\RaggedRight #1}\par}

% Celda marcable dentro de la rubrica (tabla de autoevaluacion). Misma
% sangria francesa, pero pensada para vivir dentro de una columna X.
\newcommand{\milcopcion}[1]{%
\makebox[5.2mm][l]{\raisebox{-.55mm}{\textcolor{milcNegro}{\milcmarca}}}%
\parbox[t]{\dimexpr\linewidth-5.2mm\relax}{\RaggedRight #1}}

% ── Recursos visuales de la guía ──────────────────────────────────────
% Declarados en el bloque `recursos:` del YAML y emitidos por el builder.
% \guiaFigura   → imagen rasterizada o PDF (foto, carta celeste, montaje).
% \guiaDiagrama → fuente TikZ/circuitikz (.tex) incrustada como vector:
%                 es la vía para circuitos y esquemáticos editables.
% [#1] = texto alternativo (alt) · #2 = ruta relativa al .tex · #3 = pie
% El alt llega al PDF etiquetado (/Alt) y lo lee el lector de pantalla.
\newcommand{\guiaFigura}[3][]{%
\begin{center}
\includegraphics[width=0.88\linewidth,height=0.38\textheight,keepaspectratio,alt={#1}]{#2}\\[1.5mm]
{\footnotesize\itshape\textcolor{milcNegro!75}{#3}}
\end{center}
}

\newcommand{\guiaDiagrama}[2]{%
\begin{center}
\input{#1}\\[1.5mm]
{\footnotesize\itshape\textcolor{milcNegro!75}{#2}}
\end{center}
}

\begin{document}
\thispagestyle{empty}

% ── Portada ───────────────────────────────────────────────────────────
% Antes: un rectangulo de color a pagina completa. Se veia bien en pantalla,
% pero en la fotocopiadora del colegio se comia el toner y salia gris sucio.
% Ahora la banda ocupa el tercio superior y el resto va en blanco.
\begin{tikzpicture}[remember picture,overlay]
  \path[fill=milcGradePrimary]
    (current page.north west) rectangle ([yshift=-10.6cm]current page.north east);

  \node[anchor=north west,text=milcGradeText] at ([xshift=2.0cm,yshift=-1.55cm]current page.north west)
    {\sffamily\bfseries\fontsize{12}{14}\selectfont GRADO};
  \node[anchor=north west,text=milcGradeText,opacity=.85] at ([xshift=2.0cm,yshift=-2.35cm]current page.north west)
    {\sffamily\bfseries\fontsize{8.5}{10}\selectfont SERIE GUÍAS · TECNOLOGÍA E INFORMÁTICA};
  \node[anchor=north east,text=milcGradeText,opacity=.85,align=right] at ([xshift=-2.0cm,yshift=-1.55cm]current page.north east)
    {\sffamily\bfseries\fontsize{9}{12}\selectfont I.E. Sor María Juliana\\Cartago, Valle del Cauca};

  \node[anchor=west,text=milcGradeText] (gradonum) at ([xshift=1.85cm,yshift=-7.1cm]current page.north west)
    {\sffamily\bfseries\fontsize{112}{112}\selectfont 9};
  % El circulito de grado (°) solo aplica a las guias de grado («11°»); en
  % Bebras y Semillero el numero grande es un momento/modulo, no un grado.
  % Se posiciona relativo al nodo `gradonum`, no en coordenadas absolutas.
  \draw[milcGradeText,line width=1.6pt]
    ([xshift=2.0mm,yshift=-4.5mm]gradonum.north east) circle (.15cm);

  \node[anchor=west,text=milcGradeText,text width=11.4cm,align=left] at ([xshift=1.5cm,yshift=0.5cm]gradonum.east)
    {
     {\sffamily\bfseries\fontsize{9}{11}\selectfont\MakeUppercase{Período 1 · Historia de la técnica}}\\[3mm]
     {\sffamily\bfseries\fontsize{21}{25}\selectfont\setlength{\baselineskip}{27pt}Cosecha del periodo ---\\[4pt]manifiesto del\\[4pt]técnico crítico}\\[3.5mm]
     {\sffamily\bfseries\fontsize{15}{18}\selectfont Guía 10-9-TIC}
    };
\end{tikzpicture}

\vspace*{9.4cm}
\vfill

{\sffamily\bfseries\fontsize{8}{10}\selectfont\textcolor{milcGradeDark}{LO QUE TE LLEVAS HOY}}

\begin{tcolorbox}[enhanced,colback=white,colframe=milcNegro,arc=2mm,boxrule=1.6pt,
  left=5mm,right=5mm,top=4mm,bottom=4mm,before skip=3pt,after skip=12pt,
  fontupper=\RaggedRight\fontsize{11}{15.5}\selectfont]
Tu manifiesto del técnico crítico en una página, con cinco elementos visibles: una posición en una frase, dos o tres razones que la sostienen, un ancla local del periodo, los límites que le reconoces y una propuesta que puedas empezar esta semana. Lo defiendes en voz alta de tres a cinco minutos, con notas breves y sin leer la hoja.
\end{tcolorbox}

\begin{tcolorbox}[enhanced,colback=milcGris,colframe=milcGris,arc=2mm,boxrule=0pt,
  left=5mm,right=5mm,top=3.5mm,bottom=3.5mm,after skip=12pt]
{\sffamily\bfseries\fontsize{8}{10}\selectfont\textcolor{milcNegro!55}{ESTA GUÍA ES DE}}\\[3mm]
\begin{tabularx}{\linewidth}{X p{3.2cm} p{3.2cm}}
\linead{\linewidth} & \linead{3.0cm} & \linead{3.0cm}\\[-1mm]
{\fontsize{8.5}{10}\selectfont\textcolor{milcNegro!55}{Nombre completo}} &
{\fontsize{8.5}{10}\selectfont\textcolor{milcNegro!55}{Grupo}} &
{\fontsize{8.5}{10}\selectfont\textcolor{milcNegro!55}{Fecha}}\\
\end{tabularx}
\end{tcolorbox}

\vfill

\noindent\textcolor{milcLinea}{\rule{\linewidth}{1pt}}\\[2mm]
{\sffamily\bfseries\fontsize{9.5}{12}\selectfont Autor: Dr. Álvaro Cárdenas Orozco}\\[1mm]
{\sffamily\fontsize{8.5}{11}\selectfont\textcolor{milcNegro!55}{Metodología MILC · apertura ancestral · pensamiento computacional · triángulo Dussel–estoicismo–Floridi}}

\newpage

\section*{Apertura: Cosecha P1 --- manifiesto del técnico crítico}

\begin{softbox}{milcGradeDark}{milcGradeSoft}{Saber ancestral}
En los pueblos indígenas del Cauca, las normas que rigen a todos se llaman \textbf{mandatos}, y no las escribe una autoridad sola. Salen de un congreso al que llega la gente de los resguardos a discutir durante días. El XV Congreso del CRIC se reunió en Río Blanco, Sotará, del 25 al 30 de junio de 2017. De allí salieron mandatos concretos que la organización después tiene que rendir públicamente: cuáles se cumplieron y cuáles no. Un mandato no vale por estar escrito. Vale porque quienes lo hicieron vuelven a mirarlo. Conviene decir esto con precisión: así construye su norma un pueblo que disputa su gobierno con el Estado y con actores armados. Lo que tú vas a hacer hoy es algo mucho más pequeño que aprende de eso. La \textbf{cara de exclusión}: la propia organización reporta que el cumplimiento de sus mandatos queda incompleto.\par\smallskip{\footnotesize\textcolor{milcNegro!70}{\textbf{Fuente.} Consejo Regional Indígena del Cauca. (2017). \emph{Avanza el XV congreso regional del CRIC en el resguardo de Rioblanco Sotará del pueblo indígena yanacona}.}}
\end{softbox}

\begin{softbox}{milcTurquesa}{milcGris}{Saber contemporáneo}
Un manifiesto tiene cinco elementos y ninguno sobra. \textbf{Posición}: una afirmación en una frase, que empiece por «creo que». \textbf{Razones}: dos o tres, y técnicas, no de gusto personal. \textbf{Ancla}: un caso concreto del periodo que sostenga lo que dices. \textbf{Límites}: qué no sabes y en qué tu posición es débil. \textbf{Propuesta}: algo que puedas empezar esta semana. La regla de redacción es «firme pero matizado». Sostener con fuerza y reconocer lo flojo no se contradicen: la rigidez suena a que no lo pensaste, y la tibieza a que no te importa. Y hay una prueba que no falla. Léelo en voz alta y pregúntate si puedes defenderlo sin avergonzarte y sin inflarte. Si te avergüenza, arréglalo. Si te infla, quítale adjetivos.
\end{softbox}

\begin{softbox}{milcMostaza}{milcGris}{Pregunta-puente}
\textit{Un mandato vale porque quienes lo escribieron vuelven a mirarlo. ¿Cuándo vas a volver a mirar el manifiesto que escribas hoy?}
\end{softbox}

\begin{softbox}{milcVerde}{milcGris}{Saber y saber hacer}
Al terminar podrás: (1) \textbf{analizar} tus propias respuestas del periodo y ver qué cambió en lo que pensabas; (2) \textbf{explicar} los cinco elementos de un manifiesto y por qué reconocer límites lo fortalece; (3) \textbf{crear} tu manifiesto en una página y defenderlo en voz alta ante el grupo.
\end{softbox}



\small
\begin{tabularx}{\textwidth}{>{\bfseries}p{3.0cm}Y}
\rowcolor{milcGradePrimary!12}Autor & Dr. Álvaro Cárdenas Orozco\\
DBA & Análisis crítico de la historia de la tecnología y su impacto social (MEN, T\&I)\\
Referentes & Dussel (técnica situada) · Estoicismo (téchne del cuerpo) · Floridi (infoesfera)\\
Producto final & Tu manifiesto del técnico crítico en una página, con cinco elementos visibles: una posición en una frase, dos o tres razones que la sostienen, un ancla local del periodo, los límites que le reconoces y una propuesta que puedas empezar esta semana. Lo defiendes en voz alta de tres a cinco minutos, con notas breves y sin leer la hoja.\\
\end{tabularx}
\normalsize

\puente{Nueve sesiones recorrieron la técnica desde el cuerpo hasta el chip, pasando por la balanza, la imprenta, la corriente y las tecnologías propias. Hoy cierras tomando posición. En el periodo 2 pasas al diseño editorial, y esa posición te va a servir.}

\milcestacion{milcGradeDark}{Ruta MILC: así vamos a aprender}
\begin{tabularx}{\textwidth}{>{\centering\arraybackslash}X>{\centering\arraybackslash}X>{\centering\arraybackslash}X>{\centering\arraybackslash}X}
\phasecard{1}{milcVerde}{milcGris}{Escucha\\[-1mm]\small Observo, escucho y pregunto.} &
\phasecard{2}{milcTurquesa}{milcGris}{Entender\\[-1mm]\small Sistematizo y ordeno.} &
\phasecard{3}{milcMagenta}{milcGris}{Hacer\\[-1mm]\small Construyo y aplico.} &
\phasecard{4}{milcVino}{milcGris}{Cierre\\[-1mm]\small Evalúo y reflexiono.}
\end{tabularx}


\puente{Antes de escribir vas a releerte. Las nueve preguntas puente del periodo están en tu cuaderno, y algunas las contestaste distinto de como las contestarías hoy.}

\milcestacion{milcVerde}{Estación 1 de 3 · Escucha}

\begin{softbox}{milcVerde}{milcGris}{Escena para escuchar}
\textbf{Actividad 1 · ANALIZA --- Las nueve preguntas del periodo} (15 min · individual). (1) Abre tu cuaderno y busca las nueve preguntas puente del periodo, la que cierra la apertura de cada sesión. (2) Anótalas en orden. (3) Marca las tres que te costaron más. (4) Escribe al lado de dos de ellas qué contestarías hoy, y si cambió respecto a lo que contestaste entonces. (5) Elige la que más te marcó: de ahí sale tu posición.
\end{softbox}

{\sffamily\bfseries\fontsize{8}{10}\selectfont\textcolor{milcNegro!55}{MARCA CUANDO LO HAYAS HECHO}}
\milccheckrow{Tengo las nueve preguntas puente anotadas en orden.}
\milccheckrow{Marqué las tres que me costaron más.}
\milccheckrow{Elegí la que más me marcó y sé por qué.}

\begin{infoband}{milcVerde}


\textbf{Lo que va al cuaderno (Actividad 1 · ANALIZA).} \textbf{Título:} «Actividad 1 --- Las nueve preguntas del periodo». \textbf{Formato:} la lista de las nueve con las tres difíciles marcadas, más dos respuestas de hoy comparadas con las de entonces. \textbf{Extensión:} un tercio de página. \textbf{Sabes que terminaste cuando} tienes elegida la pregunta de la que va a salir tu posición.
\end{infoband}

\puente{Ya tienes tu material. Ahora vas a ver qué separa un manifiesto de una opinión suelta, y por qué reconocer un límite lo hace más fuerte, no más débil.}

\milcestacion{milcTurquesa}{Estación 2 de 3 · Entender}

\begin{softbox}{milcTurquesa}{milcGris}{Lo que vas a entender}
\textbf{Actividad 2 · EXPLICA --- Los cinco elementos} (20 min · parejas). (1) Con tu pareja, escriban los cinco elementos con una frase propia cada uno. (2) Escriban cada uno su posición en una sola frase y léansela al otro. (3) Búsquenle a la posición del otro un límite que no haya visto. (4) Anoten por qué reconocer ese límite hace la posición más creíble y no menos.
\end{softbox}

\milcpaso{1}{milcTurquesa}{\textbf{Los cinco elementos.} Posición en una frase. Dos o tres razones técnicas. Un ancla local del periodo. Los límites reconocidos. Una propuesta que empiece esta semana.}
\milcpaso{2}{milcTurquesa}{\textbf{Firme pero matizado.} Sostener con fuerza y reconocer lo flojo no se contradicen. La rigidez suena a que no lo pensaste; la tibieza, a que no te importa. El punto medio no es blandura: es haber mirado tu propia posición por el revés.}
\milcpaso{3}{milcTurquesa}{\textbf{Las razones son técnicas, no de gusto.} «No me gusta que todo sea digital» es un gusto. «Delegar un cálculo sin decidirlo deja una capacidad sin uso» es una razón, y viene de la sesión 8.}
\milcpaso{4}{milcTurquesa}{\textbf{Lo que escribes no es un mandato.} Un mandato es un acto de gobierno de un pueblo, discutido en congreso y rendido en público. Lo tuyo es mucho más pequeño y aprende de eso: una posición propia que te comprometes a volver a mirar.}

\begin{drawbox}{milcTurquesa}{Los cinco elementos de un manifiesto}
\textbf{Posición:} una frase que empieza por «creo que». \\[1mm] \textbf{Razones:} dos o tres, técnicas. \\[1mm] \textbf{Ancla:} un caso concreto del periodo. \\[1mm] \textbf{Límites:} qué no sé y dónde soy débil. \\[1mm] \textbf{Propuesta:} algo que empiece esta semana.
\end{drawbox}

\begin{softbox}{milcMostaza}{milcGris}{Errores comunes y su corrección}
(1) \textbf{Posición de dos párrafos}: si no cabe en una frase, todavía no la tienes. (2) \textbf{Razones que son gustos}: lo que te agrada no sostiene nada. (3) \textbf{Ancla genérica}: «la tecnología ancestral» no es un ancla; el trapiche del río La Vieja sí. (4) \textbf{Sin límites}: un texto sin una sola duda se lee como propaganda. (5) \textbf{Leer la hoja al sustentar}: si la lees, sobra quien la escribió.
\end{softbox}

\begin{infoband}{milcTurquesa}


\textbf{Lo que va al cuaderno (Actividad 2 · EXPLICA).} \textbf{Título:} «Actividad 2 --- Los cinco elementos». \textbf{Formato:} los cinco con frase propia, tu posición en una sola frase y el límite que te encontró tu pareja. \textbf{Extensión:} media página. \textbf{Sabes que terminaste cuando} tienes escrito un límite que no habías visto solo.
\end{infoband}

\puente{Tienes la estructura. Toca escribir la página, ensayarla y defenderla delante del grupo.}

\milcestacion{milcMagenta}{Estación 3 de 3 · Hacer}

\begin{softbox}{milcMagenta}{milcGris}{Cómo vas a construir tu producto}
\textbf{Actividad 3 · CREA --- Tu manifiesto y su defensa} (30 min · individual). (1) Escribe tu posición en una frase, arriba de la página. (2) Añade dos o tres razones técnicas, cada una en dos renglones. (3) Añade el ancla local del periodo que la sostiene. (4) Escribe los límites que le reconoces, sin suavizarlos. (5) Cierra con una propuesta que puedas empezar esta semana. (6) Ensáyalo con cronómetro y defiéndelo ante el grupo con notas breves, sin leer la hoja.
\end{softbox}

\milcpaso{1}{milcMagenta}{\textbf{Tu entrega tiene dos piezas.} La página con los cinco elementos visibles a simple vista. Y la sustentación de tres a cinco minutos, con notas y no con la hoja.}
\milcpaso{2}{milcMagenta}{\textbf{El ancla la eliges tú.} El jaibaná que aprende acompañando, las dos medidas de Saucío, los oficios asignados de la finca cafetera, la décima que vuelve al inicio, la radio con NIT, el molde que se disuelve, las Pléyades. Escoge la que de verdad sostenga tu posición.}
\milcpaso{3}{milcMagenta}{\textbf{El párrafo de límites es el que más pesa.} Escribir qué no sabes es lo que separa una posición pensada de una consigna. Y en la sustentación es lo que hace que te crean.}
\milcpaso{4}{milcMagenta}{\textbf{La propuesta tiene fecha.} «Voy a pensar más antes de usar la calculadora» no se puede cumplir. «Del lunes al viernes calculo el vuelto de memoria» sí, y el viernes se sabe.}

\begin{drawbox}{milcMagenta}{Antes de sustentar}
\checkbox Posición en una sola frase, visible arriba. \\[1mm] \checkbox Dos o tres razones técnicas, no gustos. \\[1mm] \checkbox Un ancla local concreta del periodo. \\[1mm] \checkbox Límites reconocidos sin suavizar. \\[1mm] \checkbox Propuesta con día y acción. \\[1mm] \checkbox Ensayado con cronómetro, entre tres y cinco minutos. \\[1mm] \checkbox Notas breves, no la hoja completa.
\end{drawbox}

\begin{softbox}{milcMagenta}{milcGris}{Plantilla de trabajo}
\textbf{Posición.} \textit{Creo que la técnica debe…} \\[1mm] \textbf{Razones.} \textit{Dos o tres, técnicas.} \\[1mm] \textbf{Ancla.} \textit{El caso del periodo que la sostiene.} \\[1mm] \textbf{Límites.} \textit{Qué no sé y dónde soy débil.} \\[1mm] \textbf{Propuesta.} \textit{Qué empiezo y qué día.}
\end{softbox}

\begin{infoband}{milcMagenta}


\textbf{Lo que va al cuaderno (Actividad 3 · CREA).} \textbf{Título:} «Actividad 3 --- Tu manifiesto y su defensa». \textbf{Formato:} el manifiesto completo con los cinco elementos y las notas breves con las que lo sustentaste. \textbf{Extensión:} una página. \textbf{Sabes que terminaste cuando} lo defendiste sin leer la hoja y sin avergonzarte.
\end{infoband}

\puente{Lo que entregas hoy: una página con los cinco elementos y una sustentación de tres a cinco minutos. La prueba de calidad: alguien que no estuvo en clase entiende tu posición leyéndola.}

\milcestacion{milcMostaza}{Producto de la sesión}
\textbf{Manifiesto del técnico crítico en una página + sustentación de tres a cinco minutos}.
\begin{softbox}{milcMostaza}{milcGris}{Criterios de calidad}
(1) La posición cabe en una frase y está visible arriba de la página. (2) Las razones son técnicas y vienen del periodo, no son gustos. (3) El ancla local es un caso concreto, no una categoría. (4) Los límites están escritos sin suavizar. (5) La propuesta dice qué empiezas y qué día.
\end{softbox}

\puente{Antes de cerrar, mira el periodo desde cinco lados. Una posición que nunca se vuelve a mirar se convierte en una frase bonita y nada más.}

\milcestacion{milcVino}{Evaluación liberadora · cinco dimensiones}

\begin{softbox}{milcVino}{milcGris}{Sentido de la evaluación}
Evaluar no es solo revisar una nota. Es reconocer qué aprendí, qué cambió en mí, qué emociones manejé, cómo me relaciono con la ciudad y la comunidad, y qué le dejaré a quien viene detrás.
\end{softbox}

\textbf{Mi reflexión liberadora en cinco dimensiones.} Completa cada frase iniciada:

\small
\begin{tabularx}{\textwidth}{>{\bfseries\RaggedRight\arraybackslash}p{3.4cm}Y>{\RaggedRight\arraybackslash}p{4.6cm}}
\rowcolor{milcVino}\color{white}Dimensión & \color{white}\bfseries Mi respuesta (frame starter) & \color{white}\bfseries Algo que descubrí\\
1. Personal & Lo que más cambió en mí fue \linead{6.5cm} & \linead{4.2cm}\\
2. Emocional & La emoción más fuerte fue \linead{2.5cm} y la regulé \linead{3cm} & \linead{4.2cm}\\
3. Ciudadana & En democracia esto importa porque \linead{6cm} & \linead{4.2cm}\\
4. Local (Cartago) & En mi barrio sería útil para \linead{6.5cm} & \linead{4.2cm}\\
5. Intergeneracional & A mi abuela/abuelo le diría \linead{6.5cm} & \linead{4.2cm}\\
\end{tabularx}
\normalsize

\begin{infoband}{milcVino}
\textbf{Para la próxima sustentación.} Vuelve a esta tabla en seis meses. Verás que las respuestas cambian — no porque lo aprendido cambie, sino porque tú cambias. La evaluación liberadora es una conversación contigo a través del tiempo.
\end{infoband}


\puente{Para terminar, tres citas verificadas sobre elegir y sostener: Dussel sobre saber perder tiempo para elegir los temas que importan, Marco Aurelio sobre cambiar de parecer con gusto, Floridi sobre saber preguntar y responder.}

\milcestacion{milcGradeDark}{Triángulo de pensamiento · cierre filosófico}

\begin{softbox}{milcVino}{milcGris}{Dussel · lente del nosotros}
\textit{«Como los temas son infinitos y el tiempo corto, es necesario saber perder tiempo para elegir los temas fundamentales de la época que nos toca vivir.»}

{\footnotesize\itshape\color{milcNegro!70}--- Enrique Dussel · Filosofía de la liberación (1977), §5.9.4.1}

Los quince minutos que dedicas a releer tus nueve preguntas no son tiempo perdido: son la elección. Un manifiesto sobre lo primero que se te ocurra no dice nada de ti; uno sobre lo que de verdad te quedó dando vueltas, sí.

\textbf{De todo lo del periodo, ¿qué me siguió dando vueltas después de la clase?}
\end{softbox}
\linead{\linewidth}\\[1mm]
\linead{\linewidth}

\begin{softbox}{milcMostaza}{milcGris}{Estoicismo · Marco Aurelio · Meditaciones VI, 21 (c. 175 d.C.) · cuidado interior}
\textit{«Si alguno pudiese convencerme y hacerme patente que yo no pienso con rectitud u obro bien, con mucho gusto mudaría de dictamen y de conducta; porque yo busco la verdad, la cual jamás dañó a ninguno.»}

{\footnotesize\itshape\color{milcNegro!70}--- Marco Aurelio · Meditaciones VI, 21 (c. 175 d.C.)}

Esa frase es el párrafo de límites, dicho hace mil ochocientos años. Sostener una posición y estar dispuesto a cambiarla no se contradicen: es lo que separa pensar de repetir.

\textbf{¿Qué tendría que mostrarme alguien para que yo cambiara la posición que escribí hoy?}
\end{softbox}
\linead{\linewidth}\\[1mm]
\linead{\linewidth}

\begin{softbox}{milcTurquesa}{milcGris}{Floridi · lente de la infoesfera}
\textit{«La partida la ganarán quienes «sepan preguntar y responder» (Platón, Crátilo 390c) y, por tanto, sepan qué datos pueden ser útiles y relevantes, y merecen recogerse y cuidarse. (trad. propia)»}

{\footnotesize\itshape\color{milcNegro!70}--- Luciano Floridi · Big data and their epistemological challenge (2012)}

En la sustentación se ve. Quien preparó su posición sabe qué preguntas puede recibir y cuáles no sabe contestar todavía. Anticipar la pregunta difícil vale más que ensayar la respuesta fácil.

\textbf{¿Cuál es la pregunta que menos quiero que me hagan sobre mi manifiesto?}
\end{softbox}
\linead{\linewidth}\\[1mm]
\linead{\linewidth}

\begin{infoband}{milcNegro}
\textbf{Nota para el docente.} Las tres son citas textuales y llevan su referencia debajo. Puedes remitir a la obra en clase.
\end{infoband}

\milcestacion{milcVino}{Mi autoevaluación · marca una casilla por fila}

{\small
\setlength{\extrarowheight}{2.2pt}
\begin{tabularx}{\textwidth}{>{\RaggedRight\arraybackslash\bfseries}p{3.0cm}YYY}
\rowcolor{milcVino}\color{white}\bfseries Criterio & \color{white}\bfseries Lo logré & \color{white}\bfseries En proceso & \color{white}\bfseries Necesito apoyo\\[.6mm]
Comprensión del tema & \milcopcion{Explico las ideas con mis palabras.} & \milcopcion{Reconozco las ideas, falta precisión.} & \milcopcion{Necesito repasar lo básico.}\\[.8mm]
\rowcolor{milcGris}Pensamiento computacional & \milcopcion{Usé los pilares al armar mi producto.} & \milcopcion{Usé algunos pilares.} & \milcopcion{Necesito releer la estación 2.}\\[.8mm]
Aplicación y producto & \milcopcion{Mi producto cumple los criterios.} & \milcopcion{Está armado, con ajustes pendientes.} & \milcopcion{Aún está en construcción.}\\[.8mm]
\rowcolor{milcGris}Reflexión 5 dimensiones & \milcopcion{Respondí cada una con honestidad.} & \milcopcion{Respondí algunas con detalle.} & \milcopcion{Mis respuestas son muy generales.}\\[.8mm]
Triángulo de pensamiento & \milcopcion{Conecté las 3 voces con mi trabajo.} & \milcopcion{Conecté al menos una.} & \milcopcion{No las relaciono con mi proceso.}\\[.8mm]
\end{tabularx}}

\vspace{3mm}
\textbf{Hoy entendí que:}\\[1mm]
\linead{\linewidth}\\[1mm]
\linead{\linewidth}

\textbf{Mi compromiso para la próxima sesión es:}\\[1mm]
\linead{\linewidth}\\[1mm]
\linead{\linewidth}

\begin{softbox}{milcMostaza}{milcGris}{Cita de despedida · Marco Aurelio}
\textit{``El obstáculo en el camino se vuelve el camino. Lo que estorba al camino se vuelve camino.''}\\[1mm]
Cada error de hoy, bien observado, se convierte en pieza del camino que pisarás mañana.
\end{softbox}

\small
\begin{tabularx}{\textwidth}{>{\bfseries}p{3.6cm}Y}
\rowcolor{milcGradeDark!12}Nombre completo: & \linead{10cm}\\
Fecha: & \linead{4cm} \quad Curso/grupo: \linead{4cm}\\
Mi firma: & \linead{9cm}\\
\end{tabularx}
\normalsize

\begin{infoband}{milcGradeDark}
\textbf{Para el periodo que viene.} Guarda el manifiesto donde puedas encontrarlo. Pon una fecha dentro de un mes para volver a leerlo y anotar si cumpliste la propuesta. En el periodo 2 entras al diseño editorial, y ahí tu posición sobre la técnica se pone a prueba con piezas reales.
\end{infoband}

\end{document}
